\clearpage
\section{Anexo 8.F: Respuestas, esfuerzo y riesgo residual proyectado}
\label{sec:cierre-respuestas}
Las 35 fichas complementan 8.A y T-16. A es esfuerzo preventivo adicional total en la ventana, no costo de desarrollo ni pruebas ordinarias de la EDT. Se estima una revisión extraordinaria de excepción/coordinación/cierre por riesgo: 60\% del dueño y 40\% de CAL (ARQ cuando CAL es dueño), acumulados en toda la ventana. Para vigilancia ordinaria y aceptación pasiva A=0. C es esfuerzo de respuesta condicional total, repartido en esos perfiles y contenido en el PERT H0 de 8.B, no sumado de nuevo ni multiplicado otra vez por episodios n. Son hipótesis de dos perfiles, no compromisos ni rendimientos observados. JP debe contrastar estos esfuerzos con los paquetes existentes antes de aprobarlos; si una acción ya está presupuestada, se descuenta aquí y se recalculan sumas.

El plazo preventivo es el de 8.A.2 para las 28 dependencias, antes del hito/fecha allí indicado; para los siete externos se fija en cada ficha. A se ejecuta una vez como presupuesto acumulado y se programa a D-5 del primer hito afectado; si hay revisiones repetidas (p. ej. rotación), se reparte ese total, sin multiplicarlo por 44 meses ni por n. Las HH preventivas mínimas no sustituyen capacitación, mantenimiento, respaldo o controles obligatorios incluidos en el presupuesto ordinario.

La probabilidad posterior p1 es un juicio de planificación de ocurrencia del mismo evento y ventana de p0, condicionado a que se ejecuten los controles y se obtenga la evidencia indicada. No es frecuencia medida ni eficacia certificada; se mantiene la probabilidad de dependencias externas que LafroX no puede modificar. En historia sanitaria, retirada del planificador, energía, SII, siniestros y cambio normativo no se reduce p0 por describir una respuesta. Los otros p1 proponen menor ocurrencia por resolver temprano la dependencia especificada; esa hipótesis se invalida si no existe su acta o prueba. I1 conserva el impacto global inicial, incluyendo consecuencias sanitarias y contractuales: ni el bloqueo de salida ni una mejora de detección borran una pérdida histórica. E1 se obtiene de la banda de p1 y el factor de I1 con los mismos umbrales de 8.1.3.

El esfuerzo condicional remanente de atención es L=H0-C; no es una nueva estimación de consecuencias sanitarias ni de plazo. El esfuerzo residual esperado es p1 L. El tratamiento esperado es T=A+p1 C. El beneficio bruto en HH es B=p0 H0-p1 L. B/T compara reducción esperada de trabajo de atención con trabajo de tratamiento; el saldo neto B-T=(p0-p1)H0-A evita doble conteo. Así, costo total posterior=T+p1 L=A+p1 H0. Estos indicadores no valorizan alimentos, multas, seguridad de personas ni aceptación contractual; un retorno menor que uno se justifica por obligación o continuidad, sin eliminar el cálculo. Las respuestas contingentes no constituyen prevención adicional gratuita.

La base p0/O/M/P corresponde a las ventanas de exposición del Anexo 8.B. Los intervalos de programación de 8.H son hipótesis de distribución de HH, no cambian la exposición descrita en 8.B.2 ni convierten HH en días. D-5 y D-10 son días hábiles, salvo los minutos y turnos expresamente indicados. Los tiempos propuestos son máximos internos y se acortan cuando el SLA o el procedimiento de incidente exija actuar antes; no amplían compromisos de servicio ni autorizan indisponibilidad en ventanas críticas. Las fechas de respuesta no sustituyen los diez días hábiles contractuales de revisión/subsanación.

\begingroup
\renewcommand{\thetable}{8.F.1}
\begin{landscape}
\setlength{\textwidth}{\linewidth}
\renewcommand{\arraystretch}{1.08}
\begin{tablalafrox}{Fichas ejecutables de los 35 riesgos. Fuente: 8.A y estimaciones adicionales de planificación de S8.}{tab:cierre-fichas}{G{0.06} Y F{0.29} F{0.15} F{0.20}}{\cab{ID} & \cab{Acción preventiva y evidencia} & \cab{Contingencia y plazo desde disparador} & \cab{Perfiles / HH adicionales} & \cab{Momento y regla de programación}}
R-001 & Solicitar al Gerente Comercial, con TI y cada cadena, agenda EDI en mes 13, confirmada en mes 14, y certificación a D-5 del inicio de las cuatro semanas finales de evidencia de E2. Ante falta de respuesta, JP escala por CT; continuar pruebas independientes sin acreditar la cadena pendiente ni sustituir EDI por carga manual. Evidencia de salida: Agenda EDI y certificación de cada cadena, con pruebas de intercambio y registro de pendientes. & Al D-5 sin certificación, JP escala en un día hábil y conserva pendiente la cadena; sólo ampliar tras su certificación. & ARQ / CAL; A: 4,80/3,20; C: 9,60/6,40 HH. & Ventana de programación 13--20. A acumulado a D-5 del primer hito afectado; para repeticiones, repartir el total. Aplicar fechas de 8.A.2 y el plazo contingente de esta fila. \\
R-002 & Solicitar a la Jefa de Calidad, con TI y Bodega, evidencia de origen y decisión sobre lotes no reconstruibles; política V-07 al cierre del mes 2 y excepciones a D-5 de validar cada carga. Si no responde, JP escala por CT, se aísla el conjunto y queda pendiente su aceptación; no se inventa historia sanitaria. Evidencia de salida: Inventario de vínculos sanitarios por lote, evidencia de origen y resolución firmada de Calidad sobre cada excepción. & Plan adicional: DAT aísla el conjunto el mismo día; Calidad decide antes del corte. No cargar ni despachar lotes sin trazabilidad exigida; si no es reconstruible, elevar decisión de alcance contractual, sin inventar historia. & DAT / CAL; A: 7,20/4,80; C: 14,40/9,60 HH. & Ventana de programación 1--20. A acumulado a D-5 del primer hito afectado; para repeticiones, repartir el total. Aplicar fechas de 8.A.2 y el plazo contingente de esta fila. \\
R-003 & Solicitar al Jefe de Abastecimiento y a TI el historial de tres años y validación de V-09 al cierre del mes 3; corregir extractos a D-5 de validar la carga. Ante silencio, escalar por CT, separar períodos incompletos y continuar el perfilado disponible sin reducir la historia exigida. Evidencia de salida: Extracto de tres años perfilado y reconciliado con Abastecimiento/TI; acta V-09. & A D-5 sin extracto válido, separar períodos y pedir extracción complementaria; revalidar antes de cargar. & DAT / CAL; A: 3,60/2,40; C: 9,60/6,40 HH. & Ventana de programación 1--18. A acumulado a D-5 del primer hito afectado; para repeticiones, repartir el total. Aplicar fechas de 8.A.2 y el plazo contingente de esta fila. \\
R-004 & Solicitar a Bodega, con Operaciones y Finanzas, protocolo de conciliación en mes 2, validado en mes 3. Confirmar conteo a D-5 del corte y cerrar acta y ajustes antes de D. Si no hay respuesta, escalar por CT, repetir conciliación y bloquear el corte afectado hasta contar con evidencia aprobada. Evidencia de salida: Acta de conteo, causas de diferencias y ajustes autorizados por Bodega/Finanzas. & Bloquear corte con diferencia no explicada; DAT presenta conciliación y nueva prueba antes de la ventana autorizada. & DAT / CAL; A: 4,80/3,20; C: 12,00/8,00 HH. & Ventana de programación 1--20. A acumulado a D-5 del primer hito afectado; para repeticiones, repartir el total. Aplicar fechas de 8.A.2 y el plazo contingente de esta fila. \\
R-005 & Solicitar al Gerente de Operaciones y a Finanzas respaldos y conciliación V-08 de saldos por cliente y transportista; política en mes 2 y respuesta a D-5 de validar la carga M8 de cada ola. Ante silencio, escalar por CT y mantener M8 pendiente para esos saldos; no presumir saldo cero. Evidencia de salida: Saldos de envases aceptados por cliente y transportista, política V-08 y evidencia de M8. & Mantener saldos afectados pendientes; Finanzas resuelve antes de su carga, sin asumir saldo cero. & DAT / CAL; A: 3,60/2,40; C: 9,60/6,40 HH. & Ventana de programación 1--20. A acumulado a D-5 del primer hito afectado; para repeticiones, repartir el total. Aplicar fechas de 8.A.2 y el plazo contingente de esta fila. \\
R-006 & Solicitar a Comercial y Operaciones las horas históricas y a Finanzas los montos de rendición, con validación de V-10 al cierre del mes 1. Si no responden, escalar por CT y mantener pendiente la referencia de aceptación; la captura futura no reconstruye la historia anterior. Evidencia de salida: Fuentes históricas de horas y montos verificadas con Comercial, Operaciones y Finanzas, V-10. & Plan adicional: JP reúne áreas dentro de dos días hábiles desde bloqueo; en cinco días hábiles exige evidencia o eleva al Comité Ejecutivo la condición contractual. No sustituir historia por mediciones futuras ni aceptar resultado sin referencia. & IMP / CAL; A: 4,80/3,20; C: 9,60/6,40 HH. & Ventana de programación 1--20. A acumulado a D-5 del primer hito afectado; para repeticiones, repartir el total. Aplicar fechas de 8.A.2 y el plazo contingente de esta fila. \\
R-007 & Solicitar a Finanzas y TI el universo de deuda vigente, con respuesta inicial al cierre del mes 3 y conciliación a D-5 de validar la migración. Ante silencio, escalar por CT y mantener la cartera incompleta pendiente; preparar extracción complementaria sin omitir deuda obligatoria. Evidencia de salida: Universo de deuda vigente conciliado con Finanzas/TI y extracción complementaria comprobada. & No dar por migrada cartera incompleta; escalar a D-5 de cada carga y repetir conciliación. & DAT / CAL; A: 3,60/2,40; C: 9,60/6,40 HH. & Ventana de programación 1--18. A acumulado a D-5 del primer hito afectado; para repeticiones, repartir el total. Aplicar fechas de 8.A.2 y el plazo contingente de esta fila. \\
R-008 & Mantener trazabilidad de trámites pendientes, consulta y reintentos idempotentes hasta recibir confirmación del SII. Ante demora externa, conservar el estado pendiente y escalar según el procedimiento acordado; el silencio no equivale a aprobación. Evidencia de salida: Cola tributaria con correlación, consulta y reintento idempotente; confirmación externa conservada. & SRE registra pendiente al detectar demora, consulta dentro de 15 min y revisa cada 30 min; si afecta salida, avisa de inmediato a Finanzas/CT. No interpretar silencio como aprobación. & SRE / CAL; A: 0,00/0,00; C: 4,80/3,20 HH. & Ventana de programación 13--56. A acumulado a D-5 del primer hito afectado; para repeticiones, repartir el total. Aplicar fechas de 8.A.2 y el plazo contingente de esta fila. \\
R-009 & Solicitar a la Jefa de Calidad la matriz térmica versionada V-11 durante mes 1 y su aprobación al cierre del mes 2, antes de certificar. Si no responde, escalar por CT y mantener el tratamiento preventivo de lecturas fuera de ficha; no cerrar marcha blanca ni aprobar parámetros por el CLIENTE. Evidencia de salida: Matriz térmica V-11 versionada y firmada por Calidad por familia antes de certificación. & Sin parámetros aprobados, bloquear certificación del alcance afectado y escalar a D-5; no relajar bloqueo térmico. & CAL / ARQ; A: 3,60/2,40; C: 7,20/4,80 HH. & Ventana de programación 1--20. A acumulado a D-5 del primer hito afectado; para repeticiones, repartir el total. Aplicar fechas de 8.A.2 y el plazo contingente de esta fila. \\
R-010 & Solicitar al Jefe de TI, con Operaciones, calendario de equipos en mes 1, entrega confirmada a D-5 y recepción técnica antes de la instalación o prueba D. Ante silencio o faltante, escalar por CT, continuar trabajo independiente y bloquear las actividades sin equipo. Evidencia de salida: Recepción técnica por equipo y sitio, con conformidad de TI/Operaciones antes de la ola. & A D-5 sin equipos, JP decide reprogramación de trabajo independiente en un día hábil; no certificar con sustitutos no validados. & IMP / CAL; A: 2,40/1,60; C: 7,20/4,80 HH. & Ventana de programación 1--20. A acumulado a D-5 del primer hito afectado; para repeticiones, repartir el total. Aplicar fechas de 8.A.2 y el plazo contingente de esta fila. \\
R-011 & Solicitar al Gerente de Operaciones y al representante de cada transportista el acuerdo y la nómina desde mes 2, firmados a D-5 de su incorporación a la ola. Ante silencio, escalar por CT y no habilitar a esa empresa; conservar su cobertura pendiente sin retirarla del alcance. Evidencia de salida: Acuerdo operativo firmado y nómina de cada transportista antes de incorporar la ola. & A D-5 sin firma, escalar por CT el mismo día; no habilitar empresa pendiente ni reducir cobertura obligatoria. & IMP / CAL; A: 2,40/1,60; C: 4,80/3,20 HH. & Ventana de programación 2--20. A acumulado a D-5 del primer hito afectado; para repeticiones, repartir el total. Aplicar fechas de 8.A.2 y el plazo contingente de esta fila. \\
R-012 & Solicitar a Operaciones y al propietario permiso de instalación por vehículo con D-10 de la visita y obtenerlo a D-5. Ante silencio, escalar por CT, suspender la intervención y proponer otra visita o vehículo autorizado; mantener pendiente la evidencia térmica. Evidencia de salida: Autorización escrita por vehículo y visita; propietario y Operaciones confirman disponibilidad. & A D-5 sin permiso, suspender visita y acordar alternativa autorizada antes de instalar. & IMP / CAL; A: 2,40/1,60; C: 7,20/4,80 HH. & Ventana de programación 1--20. A acumulado a D-5 del primer hito afectado; para repeticiones, repartir el total. Aplicar fechas de 8.A.2 y el plazo contingente de esta fila. \\
R-013 & Solicitar a Comercial y Operaciones, con Bodega por turno, responsables y agenda de acompañamiento a D-5 del inicio de cada ola, convocando a D-10. Si no responden o persiste el registro externo, escalar por CT, focalizar acompañamiento y mantener pendiente la evidencia de adopción. Evidencia de salida: Transacciones de muestra registradas en el flujo nuevo y cobertura certificada por turno/ola. & Si persiste flujo paralelo, IMP abre refuerzo en el siguiente turno y mide diariamente; sin evidencia de adopción no cerrar ola. & IMP / CAL; A: 4,80/3,20; C: 12,00/8,00 HH. & Ventana de programación 13--56. A acumulado a D-5 del primer hito afectado; para repeticiones, repartir el total. Aplicar fechas de 8.A.2 y el plazo contingente de esta fila. \\
R-014 & Solicitar al Gerente de Operaciones y al planificador confirmación de retiro y disponibilidad V-14 al cierre del mes 1; transferir y validar reglas con el ayudante en meses 1--3. Probar las rutas y excepciones sin el planificador en la marcha blanca de meses 13--15 (S7, 7.2.1). Si falta respuesta o se retira antes, escalar por CT y reconstruir con ayudante e historia verificable, sin dar por validadas reglas pendientes. Evidencia de salida: V-14, reglas y excepciones transferidas al ayudante y prueba en 13--15 sin planificador. & Retiro confirmado no cambia p0 del fallo de transferencia. Si faltan reglas, IMP inicia reconstrucción en dos días hábiles y valida antes de cerrar el piloto. & IMP / CAL; A: 3,60/2,40; C: 9,60/6,40 HH. & Ventana de programación 1--15. A acumulado a D-5 del primer hito afectado; para repeticiones, repartir el total. Aplicar fechas de 8.A.2 y el plazo contingente de esta fila. \\
R-015 & Solicitar a Operaciones, Bodega, Comercial y TI, según perfil, nómina y liberación horaria a D-5 de cada sesión convocada a D-10. Ante silencio o ausencia, escalar por CT y reprogramar evaluación; usar sólo suplentes certificados y no acreditar preparación ni cierre de ola sin evidencia. Evidencia de salida: Lista de convocados, asistencia, evaluación y suplentes certificados por sesión. & Ausencia: avisar a jefatura el mismo día y fijar recuperación dentro de cinco días hábiles, antes de trabajo autónomo. & IMP / CAL; A: 2,40/1,60; C: 7,20/4,80 HH. & Ventana de programación 1--56. A acumulado a D-5 del primer hito afectado; para repeticiones, repartir el total. Aplicar fechas de 8.A.2 y el plazo contingente de esta fila. \\
R-016 & Solicitar a CT área decisora y suplente, con agenda al cierre del mes 1; cursar insumos a D-10 y pedir respuesta a D-5. Las aceptaciones formales conservan diez días hábiles. Para promover la región de contingencia, avisar al suplente a los 15 min sin respuesta del titular y obtener decisión dentro de 30 min (S4, 4.3.2.5); no esperar el plazo ordinario. Mantener bloqueada la acción sin autorización; el silencio no permite cambios. Evidencia de salida: Matriz de decisión con titular y suplente; recepción de insumos y resolución trazables. & D-5 sin decisión: escalar un día hábil después; promoción regional: suplente a 15 min y decisión dentro de 30 min, sin actuar por silencio. & JP / CAL; A: 2,40/1,60; C: 7,20/4,80 HH. & Ventana de programación 1--56. A acumulado a D-5 del primer hito afectado; para repeticiones, repartir el total. Aplicar fechas de 8.A.2 y el plazo contingente de esta fila. \\
R-017 & Mitigar. Verificar respaldos antes de aceptar la ola de Concepción; no considerar disponible un enlace sólo por estar ofertado. Mantener y ensayar los enlaces ya ofertados: en Talca y Concepción, fibra y LTE con Starlink en espera caliente; en cross-docking, Starlink con LTE Entel/Claro y conmutación menor a 30 s. En Los Ángeles de madrugada, probar pérdida del camino Starlink restante y continuidad local de su ventana de 3 h. Coordinar restitución y conciliación sin duplicar el mismo corte como falla de SII, POS o mapas. Referencia: S4, 4.2.5.3. Evidencia de salida: Recepción de respaldos por sitio y ensayo de pérdida del último camino, incluida Concepción y Los Ángeles. & SRE activa continuidad al corte y avisa a TI/CT en 15 min; comprobar conmutación menor a 30 s, ventana local y restitución antes de normalizar. & SRE / CAL; A: 2,40/1,60; C: 7,20/4,80 HH. & Ventana de programación 1--56. A acumulado a D-5 del primer hito afectado; para repeticiones, repartir el total. Aplicar fechas de 8.A.2 y el plazo contingente de esta fila. \\
R-018 & Verificar autonomía, alarmas y procedimiento de recuperación de los componentes de LafroX, coordinando las dependencias de energía con el CLIENTE/proveedor. Si se excede la autonomía, activar recuperación y registrar la indisponibilidad e impacto. Evidencia de salida: Autonomía y transferencia verificadas por componente; coordinación de energía con TI/Operaciones. & Al corte, SRE mide autonomía restante y alerta en 15 min; activar recuperación antes de agotarla y bloquear operación insegura. & SRE / CAL; A: 0,00/0,00; C: 4,80/3,20 HH. & Ventana de programación 13--56. A acumulado a D-5 del primer hito afectado; para repeticiones, repartir el total. Aplicar fechas de 8.A.2 y el plazo contingente de esta fila. \\
R-019 & Solicitar por CT, con TI y la autoridad del CLIENTE encauzada por Gerencia General, aprobación escrita de residencia y transferencia en mes 1, resuelta al cierre del mes 2 o antes de habilitar datos, lo que ocurra primero. Ante silencio, escalar y preparar alternativa de arquitectura; mantener deshabilitada la transferencia no aprobada. Evidencia de salida: Aprobación escrita de residencia/transferencia para regiones propuestas antes de habilitar datos. & Sin aprobación al mes 2, JP escala en un día hábil y ARQ presenta alternativa en cinco días hábiles; no transferir datos sin autorización. & ARQ / CAL; A: 2,40/1,60; C: 7,20/4,80 HH. & Ventana de programación 1--2. A acumulado a D-5 del primer hito afectado; para repeticiones, repartir el total. Aplicar fechas de 8.A.2 y el plazo contingente de esta fila. \\
R-020 & Solicitar a CT, Operaciones y TI flujo y suplentes al cierre del mes 1; ante incidente, pedir autorización escrita antes del cambio. Sin respuesta, escalar en una hora o antes si lo exige el plazo de atención, manteniendo contención autorizada y estado seguro. El silencio no permite intervenir durante congelamiento. Evidencia de salida: Flujo de cambios urgentes y suplentes aprobados; autorización escrita por intervención. & Durante incidente escalar en una hora o antes si exige el servicio; contener sin intervención prohibida y conservar congelamientos. & SEG / CAL; A: 2,40/1,60; C: 7,20/4,80 HH. & Ventana de programación 1--56. A acumulado a D-5 del primer hito afectado; para repeticiones, repartir el total. Aplicar fechas de 8.A.2 y el plazo contingente de esta fila. \\
R-021 & Detener de forma segura el recorrido afectado, registrar los casos de prueba incompletos y coordinar con el CLIENTE/transportista una nueva ventana antes de repetir la validación. Evidencia de salida: Ventana y recorrido de prueba confirmados por Operaciones/transportista antes de salir. & Al aviso de falla, suspender prueba, notificar el mismo día y acordar recorrido dentro de dos días hábiles; repetir sólo con vehículo seguro. & IMP / CAL; A: 0,00/0,00; C: 4,80/3,20 HH. & Ventana de programación 1--20. A acumulado a D-5 del primer hito afectado; para repeticiones, repartir el total. Aplicar fechas de 8.A.2 y el plazo contingente de esta fila. \\
R-022 & Priorizar la atención del accidente conforme a los procedimientos del transportista y del CLIENTE; suspender la prueba, registrar su impacto y reprogramar sólo cuando exista autorización para continuar. Evidencia de salida: Autorización de recorrido y protocolo de seguridad del CLIENTE/transportista. & Suspender de inmediato ante siniestro; coordinar el mismo día y fijar nueva ventana dentro de dos días hábiles, condicionada a autorización segura. & IMP / CAL; A: 0,00/0,00; C: 4,80/3,20 HH. & Ventana de programación 1--20. A acumulado a D-5 del primer hito afectado; para repeticiones, repartir el total. Aplicar fechas de 8.A.2 y el plazo contingente de esta fila. \\
R-023 & Coordinar con el CLIENTE y el transportista la respuesta al asalto, robo o incendio, proteger a las personas y registrar qué evidencia y pruebas se perdieron. Acordar una nueva ejecución segura antes de continuar la validación. Evidencia de salida: Custodia y seguridad del transporte corresponden al CLIENTE/transportista; evidencia de prueba preservada. & Interrumpir ante asalto, avisar el mismo día y acordar reejecución dentro de dos días hábiles; no reanudar sin habilitación segura. & IMP / CAL; A: 0,00/0,00; C: 4,80/3,20 HH. & Ventana de programación 1--20. A acumulado a D-5 del primer hito afectado; para repeticiones, repartir el total. Aplicar fechas de 8.A.2 y el plazo contingente de esta fila. \\
R-024 & Solicitar a Operaciones y al transportista nómina y suplentes a D-5 de incorporar la empresa y confirmar conductor, vehículo y viaje antes de cada salida. Si falta respuesta o llega un sustituto, escalar, verificar en andén y solicitar conductor certificado; no habilitar uso no autorizado ni acreditar cobertura incompleta. Evidencia de salida: Identidad, autorización, capacitación y vínculo persona--vehículo--viaje verificados en andén. & Sustituto no certificado: bloquear acceso al detectar y solicitar otro conductor ese mismo turno; verificar antes de salida. & IMP / CAL; A: 2,40/1,60; C: 7,20/4,80 HH. & Ventana de programación 2--56. A acumulado a D-5 del primer hito afectado; para repeticiones, repartir el total. Aplicar fechas de 8.A.2 y el plazo contingente de esta fila. \\
R-025 & Solicitar por CT a la gerencia peticionaria solicitud escrita y decisión del cambio antes de su primera actividad dependiente; registrar y remitir el mismo día hábil. Ante silencio, mantener línea base y continuar trabajo independiente; si amenaza un hito, escalar al Comité Ejecutivo sin iniciar desarrollo no aprobado. Evidencia de salida: Solicitud y análisis de impacto aprobados antes de iniciar trabajo fuera de alcance. & Registrar solicitud el mismo día; JP responde o escala en dos días hábiles y mantiene línea base hasta decisión formal. & JP / CAL; A: 2,40/1,60; C: 7,20/4,80 HH. & Ventana de programación 1--56. A acumulado a D-5 del primer hito afectado; para repeticiones, repartir el total. Aplicar fechas de 8.A.2 y el plazo contingente de esta fila. \\
R-026 & Revisar cambios normativos aplicables, determinar su efecto con las áreas competentes y tramitar ajustes de controles, retención o eliminación mediante decisión y control de cambios documentados, sin posponer una obligación legal aplicable. Evidencia de salida: Seguimiento de cambios aplicables con asesoría competente; decisión contractual documentada. & ARQ registra aviso el mismo día hábil, solicita interpretación dentro de dos y presenta impacto dentro de cinco tras recibirla; no inventar multas o reformas. & ARQ / CAL; A: 0,00/0,00; C: 9,60/6,40 HH. & Ventana de programación 1--56. A acumulado a D-5 del primer hito afectado; para repeticiones, repartir el total. Aplicar fechas de 8.A.2 y el plazo contingente de esta fila. \\
R-027 & Solicitar a Finanzas, con Comercial y CT, reglas de 5.4.3 aprobadas al cierre del mes 20 y tres liquidaciones en sombra reproducidas en meses 21--23 antes del pago variable del mes 24. Ante silencio, no activar el tramo; una disputa se escala y suspende sólo el devengo correspondiente según contrato, manteniendo el servicio. Evidencia de salida: Reglas de costo de servir aprobadas al mes 20 y tres liquidaciones en sombra reproducidas en 21--23. & Antes del mes 24 sin acuerdo, no activar tramo variable; Finanzas/CT reciben disputa en un día hábil y mantienen servicio mientras resuelven. & JP / CAL; A: 3,60/2,40; C: 7,20/4,80 HH. & Ventana de programación 20--56. A acumulado a D-5 del primer hito afectado; para repeticiones, repartir el total. Aplicar fechas de 8.A.2 y el plazo contingente de esta fila. \\
R-028 & Mitigar y escalar. JP coordina con Operaciones, TI y sindicato el acuerdo 5.4.2 antes de la ola: convocar a D-10 y resolver a D-5. Explicar finalidades, acceso y controles; mantener excluidas las cámaras en cabina y no usar GPS para jornada sin acuerdo previo. Ante desacuerdo, escalar por CT, mantener sólo el uso operacional autorizado y registrar la cobertura de la ola pendiente. Evidencia de salida: Acuerdo 5.4.2 con finalidades, accesos y controles comprendidos por Operaciones/TI/sindicato. & D-5 sin acuerdo: JP escala el mismo día y conserva sólo uso operacional autorizado; no cámaras en cabina ni GPS de jornada no acordado. & JP / CAL; A: 3,60/2,40; C: 7,20/4,80 HH. & Ventana de programación 1--20. A acumulado a D-5 del primer hito afectado; para repeticiones, repartir el total. Aplicar fechas de 8.A.2 y el plazo contingente de esta fila. \\
R-029 & Mitigar. Mantener 7.1.2 en meses 13--56, turno 22:00--06:00, tutor y evaluación antes del trabajo autónomo, con aprendizaje máximo de dos horas. Bodega y Operaciones entregan altas por turno antes de asignarlas. IMP conduce hasta mes 21 y SRE desde 22, familia IMP. Ante sobrecarga, usar personal certificado, reforzar la inducción y escalar por CT; no duplicar las HH ordinarias de 7.1.2. Evidencia de salida: Altas por turno, tutor, evaluación antes de autonomía y aprendizaje dentro de dos horas según 7.1.2. & Plan adicional: Bodega asigna cobertura certificada antes del turno; IMP/SRE revisa cada alta diariamente. Si supera tutoría, no asignar trabajo autónomo; presentar refuerzo en dos días hábiles y revalidar antes de habilitar. & IMP / CAL; A: 7,20/4,80; C: 19,20/12,80 HH. & Ventana de programación 13--56. A acumulado a D-5 del primer hito afectado; para repeticiones, repartir el total. Aplicar fechas de 8.A.2 y el plazo contingente de esta fila. \\
R-030 & Mitigar y escalar. ARQ solicita a TI acceso, diccionario y condiciones de extracción en mes 1 (1.2.3, 3.3.2 y V-16); DAT verifica la migración del WMS. Probar temprano sin modificar el ERP. Si falta una interfaz, evaluar intercambio de archivos controlado, sujeto a demostrar los requisitos; elevar por CT el bloqueo y sus impactos. La alternativa no acredita por sí sola latencia, integridad ni emisión tributaria. Evidencia de salida: Contrato de extracción e intercambio confirmado con TI; ensayo temprano sin modificar ERP y sustitución progresiva de WMS. & Plan adicional: ARQ eleva brecha al finalizar semana 2 del mes 1; TI entrega acceso o alternativa verificable antes de integración. Ensayar archivo controlado si satisface requisitos; bloquear certificación si persiste interfaz no verificable. & ARQ / CAL; A: 7,20/4,80; C: 14,40/9,60 HH. & Ventana de programación 1--18. A acumulado a D-5 del primer hito afectado; para repeticiones, repartir el total. Aplicar fechas de 8.A.2 y el plazo contingente de esta fila. \\
R-031 & Mitigar y escalar. JP acuerda con Operaciones y CT disponibilidad y continuidad antes de septiembre; cerrar habilitaciones en agosto según S7. En meses 19--20 mantener evidencia diaria de las seis condiciones de marcha blanca. Un defecto durante 1--25 de septiembre no autoriza despliegue: activar continuidad aprobada y R-020 si requiere autorización. Si persiste una condición incumplida, no declarar cierre; elevar plan de recuperación sin mover hitos por cuenta propia. Evidencia de salida: Disponibilidad y continuidad acordadas antes de septiembre; seguimiento diario de las seis condiciones de cierre. & Plan adicional: cerrar habilitaciones en mes 19 antes del congelamiento; JP revisa incidentes diariamente en 19--20. Escalar el mismo día una condición incumplida y acordar recuperación en dos días hábiles, sin intervención prohibida ni aceptación ficticia. & JP / CAL; A: 4,80/3,20; C: 9,60/6,40 HH. & Ventana de programación 19--20. A acumulado a D-5 del primer hito afectado; para repeticiones, repartir el total. Aplicar fechas de 8.A.2 y el plazo contingente de esta fila. \\
R-032 & Mitigar. ARQ mantiene inventario de versiones y revisa soporte y compatibilidad trimestralmente desde mes 1, con DES y SRE. Preparar y probar adecuación antes del límite verificado de soporte, dentro de ventanas autorizadas. Si el cambio excede la actualización ordinaria 8.2.5, registrar el esfuerzo extraordinario y elevar impactos por CT; usar reversión ensayada y soporte permitido, sin sostener una versión insegura por conveniencia. Evidencia de salida: Inventario de versiones y soporte, revisión trimestral y compatibilidad probada antes del límite verificado. & Cambio excepcional: ARQ registra aviso en un día hábil y presenta transición en cinco; programar antes de fin de soporte, conservando ventanas y pruebas. & ARQ / CAL; A: 3,60/2,40; C: 12,00/8,00 HH. & Ventana de programación 1--56. A acumulado a D-5 del primer hito afectado; para repeticiones, repartir el total. Aplicar fechas de 8.A.2 y el plazo contingente de esta fila. \\
R-033 & Mitigar. ARQ verifica con TI propiedad de cuentas, acceso a datos, claves, exportaciones e infraestructura como código en la recepción inicial y en cada actualización anual del plan 8.1.4. Probar una restauración utilizable, no sólo descargar archivos. Ante restricción, documentar brecha, mantener continuidad autorizada y escalar por CT alternativas y decisión contractual; no declarar reversibilidad completa ni cambiar de proveedor sin prueba y aprobación. Evidencia de salida: Cuentas, claves, exportación y restauración utilizable verificadas en recepción y revisión anual de salida. & Restricción efectiva: ARQ escala en un día hábil y presenta alternativas en cinco; no sustituir proveedor sin prueba/aprobación ni declarar reversibilidad por una descarga. & ARQ / CAL; A: 2,40/1,60; C: 14,40/9,60 HH. & Ventana de programación 1--56. A acumulado a D-5 del primer hito afectado; para repeticiones, repartir el total. Aplicar fechas de 8.A.2 y el plazo contingente de esta fila. \\
R-034 & Mitigar y escalar. ARQ/SRE acuerdan con TI y Finanzas restitución del ERP y ensayan AL-DTE-01/AL-DR-01 antes de H5/H10, incluida pérdida de sala. Mantener sólo salidas con guía preemitida aún válida y bloquear las que requieren nueva guía. Recuperar WMS en nube no restaura el ERP ni exceptúa el RTO requerido; registrar y escalar por CT cualquier incumplimiento, sin segundo emisor ni aprobación tributaria por silencio. Evidencia de salida: Procedimiento ERP con TI/Finanzas probado en AL-DTE-01/AL-DR-01 antes de H5/H10. & Al fallo ERP avisar de inmediato a TI/Finanzas, mantener sólo guías preemitidas válidas y bloquear nuevas salidas sin guía. Activar restitución con RTO 4 h/RPO 15 min exigidos; escalar a los 30 min sin recuperación encaminada. & ARQ / CAL; A: 3,60/2,40; C: 12,00/8,00 HH. & Ventana de programación 13--56. A acumulado a D-5 del primer hito afectado; para repeticiones, repartir el total. Aplicar fechas de 8.A.2 y el plazo contingente de esta fila. \\
R-035 & Mitigar y escalar. IMP coordina con Operaciones, transportista y conductor demostración y acompañamiento antes de su salida, dentro de 7.1.3. Confirmar disponibilidad a D-5 de la ola y aceptación práctica en andén. Ante rechazo, solicitar a la empresa conductor certificado y dispuesto, mantener acompañamiento y elevar por CT la cobertura pendiente; no atribuir a LafroX facultades laborales ni declarar completa la evidencia faltante. Evidencia de salida: Demostración práctica en andén e incorporación acordada con transportista y conductor habilitado. & Ante rechazo actuar antes de salida: IMP solicita conductor certificado/dispuesto ese turno y registra cobertura pendiente; refuerzo antes de siguiente asignación. & IMP / CAL; A: 3,60/2,40; C: 9,60/6,40 HH. & Ventana de programación 2--56. A acumulado a D-5 del primer hito afectado; para repeticiones, repartir el total. Aplicar fechas de 8.A.2 y el plazo contingente de esta fila. \\
\end{tablalafrox}
\end{landscape}
\endgroup

\begingroup
\renewcommand{\thetable}{8.F.2}
\begin{landscape}
\setlength{\textwidth}{\linewidth}
\renewcommand{\arraystretch}{1.08}
\begin{tablalafrox}{Riesgo residual proyectado, condicionado a evidencia; no resultado observado. Fuente: hipótesis por riesgo y escalas de S8.}{tab:cierre-residual}{G{0.06} F{0.13} F{0.17} F{0.16} Y}{\cab{ID} & \cab{p1 / nivel / factor} & \cab{I1 / factor; E1 y zona} & \cab{Esfuerzo residual} & \cab{Fundamento, evidencia y decisión posterior}}
R-001 & 25,00\% / 2 / 0,30 & 4 / 0,40; 0,12 (Media) & L=88,00 HH; p1L=22,00 HH. & Hipótesis p1=25,00\% por resolución temprana de la dependencia: Agenda EDI y certificación de cada cadena, con pruebas de intercambio y registro de pendientes. Revisar al verificar evidencia y en comité quincenal; si falla la hipótesis, conservar la valoración inicial y activar 8.F.1. \\
R-002 & 60,00\% / 4 / 0,70 & 5 / 0,80; 0,56 (Crítica) & L=170,67 HH; p1L=102,40 HH. & No se reduce p0: Inventario de vínculos sanitarios por lote, evidencia de origen y resolución firmada de Calidad sobre cada excepción. Plan adicional en 8.F.1; mantener bloqueo/escalamiento y revisar diariamente hasta cerrar evidencia. \\
R-003 & 30,00\% / 2 / 0,30 & 3 / 0,20; 0,06 (Baja) & L=112,00 HH; p1L=33,60 HH. & Hipótesis p1=30,00\% por resolución temprana de la dependencia: Extracto de tres años perfilado y reconciliado con Abastecimiento/TI; acta V-09. Revisar al verificar evidencia y en comité quincenal; si falla la hipótesis, conservar la valoración inicial y activar 8.F.1. \\
R-004 & 30,00\% / 2 / 0,30 & 4 / 0,40; 0,12 (Media) & L=90,67 HH; p1L=27,20 HH. & Hipótesis p1=30,00\% por resolución temprana de la dependencia: Acta de conteo, causas de diferencias y ajustes autorizados por Bodega/Finanzas. Revisar al verificar evidencia y en comité quincenal; si falla la hipótesis, conservar la valoración inicial y activar 8.F.1. \\
R-005 & 30,00\% / 2 / 0,30 & 4 / 0,40; 0,12 (Media) & L=80,00 HH; p1L=24,00 HH. & Hipótesis p1=30,00\% por resolución temprana de la dependencia: Saldos de envases aceptados por cliente y transportista, política V-08 y evidencia de M8. Revisar al verificar evidencia y en comité quincenal; si falla la hipótesis, conservar la valoración inicial y activar 8.F.1. \\
R-006 & 35,00\% / 3 / 0,50 & 4 / 0,40; 0,20 (Alta) & L=85,33 HH; p1L=29,87 HH. & Hipótesis p1=35,00\% por resolución temprana de la dependencia: Fuentes históricas de horas y montos verificadas con Comercial, Operaciones y Finanzas, V-10. Plan adicional en 8.F.1; mantener bloqueo/escalamiento y revisar diariamente hasta cerrar evidencia. \\
R-007 & 25,00\% / 2 / 0,30 & 3 / 0,20; 0,06 (Baja) & L=65,33 HH; p1L=16,33 HH. & Hipótesis p1=25,00\% por resolución temprana de la dependencia: Universo de deuda vigente conciliado con Finanzas/TI y extracción complementaria comprobada. Revisar al verificar evidencia y en comité quincenal; si falla la hipótesis, conservar la valoración inicial y activar 8.F.1. \\
R-008 & 10,00\% / 1 / 0,10 & 3 / 0,20; 0,02 (Baja) & L=26,00 HH; p1L=2,60 HH. & No se reduce p0: Cola tributaria con correlación, consulta y reintento idempotente; confirmación externa conservada. Revisar al verificar evidencia y en comité quincenal; si falla la hipótesis, conservar la valoración inicial y activar 8.F.1. \\
R-009 & 20,00\% / 2 / 0,30 & 4 / 0,40; 0,12 (Media) & L=65,33 HH; p1L=13,07 HH. & Hipótesis p1=20,00\% por resolución temprana de la dependencia: Matriz térmica V-11 versionada y firmada por Calidad por familia antes de certificación. Revisar al verificar evidencia y en comité quincenal; si falla la hipótesis, conservar la valoración inicial y activar 8.F.1. \\
R-010 & 20,00\% / 2 / 0,30 & 4 / 0,40; 0,12 (Media) & L=56,00 HH; p1L=11,20 HH. & Hipótesis p1=20,00\% por resolución temprana de la dependencia: Recepción técnica por equipo y sitio, con conformidad de TI/Operaciones antes de la ola. Revisar al verificar evidencia y en comité quincenal; si falla la hipótesis, conservar la valoración inicial y activar 8.F.1. \\
R-011 & 25,00\% / 2 / 0,30 & 4 / 0,40; 0,12 (Media) & L=35,33 HH; p1L=8,83 HH. & Hipótesis p1=25,00\% por resolución temprana de la dependencia: Acuerdo operativo firmado y nómina de cada transportista antes de incorporar la ola. Revisar al verificar evidencia y en comité quincenal; si falla la hipótesis, conservar la valoración inicial y activar 8.F.1. \\
R-012 & 15,00\% / 2 / 0,30 & 4 / 0,40; 0,12 (Media) & L=47,33 HH; p1L=7,10 HH. & Hipótesis p1=15,00\% por resolución temprana de la dependencia: Autorización escrita por vehículo y visita; propietario y Operaciones confirman disponibilidad. Revisar al verificar evidencia y en comité quincenal; si falla la hipótesis, conservar la valoración inicial y activar 8.F.1. \\
R-013 & 30,00\% / 2 / 0,30 & 4 / 0,40; 0,12 (Media) & L=108,00 HH; p1L=32,40 HH. & Hipótesis p1=30,00\% por resolución temprana de la dependencia: Transacciones de muestra registradas en el flujo nuevo y cobertura certificada por turno/ola. Revisar al verificar evidencia y en comité quincenal; si falla la hipótesis, conservar la valoración inicial y activar 8.F.1. \\
R-014 & 20,00\% / 2 / 0,30 & 4 / 0,40; 0,12 (Media) & L=86,67 HH; p1L=17,33 HH. & No se reduce p0: V-14, reglas y excepciones transferidas al ayudante y prueba en 13--15 sin planificador. Revisar al verificar evidencia y en comité quincenal; si falla la hipótesis, conservar la valoración inicial y activar 8.F.1. \\
R-015 & 20,00\% / 2 / 0,30 & 4 / 0,40; 0,12 (Media) & L=73,33 HH; p1L=14,67 HH. & Hipótesis p1=20,00\% por resolución temprana de la dependencia: Lista de convocados, asistencia, evaluación y suplentes certificados por sesión. Revisar al verificar evidencia y en comité quincenal; si falla la hipótesis, conservar la valoración inicial y activar 8.F.1. \\
R-016 & 25,00\% / 2 / 0,30 & 4 / 0,40; 0,12 (Media) & L=73,33 HH; p1L=18,33 HH. & Hipótesis p1=25,00\% por resolución temprana de la dependencia: Matriz de decisión con titular y suplente; recepción de insumos y resolución trazables. Revisar al verificar evidencia y en comité quincenal; si falla la hipótesis, conservar la valoración inicial y activar 8.F.1. \\
R-017 & 20,00\% / 2 / 0,30 & 4 / 0,40; 0,12 (Media) & L=48,00 HH; p1L=9,60 HH. & Hipótesis p1=20,00\% por resolución temprana de la dependencia: Recepción de respaldos por sitio y ensayo de pérdida del último camino, incluida Concepción y Los Ángeles. Revisar al verificar evidencia y en comité quincenal; si falla la hipótesis, conservar la valoración inicial y activar 8.F.1. \\
R-018 & 5,00\% / 1 / 0,10 & 4 / 0,40; 0,04 (Baja) & L=20,67 HH; p1L=1,03 HH. & No se reduce p0: Autonomía y transferencia verificadas por componente; coordinación de energía con TI/Operaciones. Revisar al verificar evidencia y en comité quincenal; si falla la hipótesis, conservar la valoración inicial y activar 8.F.1. \\
R-019 & 20,00\% / 2 / 0,30 & 4 / 0,40; 0,12 (Media) & L=57,33 HH; p1L=11,47 HH. & Hipótesis p1=20,00\% por resolución temprana de la dependencia: Aprobación escrita de residencia/transferencia para regiones propuestas antes de habilitar datos. Revisar al verificar evidencia y en comité quincenal; si falla la hipótesis, conservar la valoración inicial y activar 8.F.1. \\
R-020 & 10,00\% / 1 / 0,10 & 4 / 0,40; 0,04 (Baja) & L=30,67 HH; p1L=3,07 HH. & Hipótesis p1=10,00\% por resolución temprana de la dependencia: Flujo de cambios urgentes y suplentes aprobados; autorización escrita por intervención. Revisar al verificar evidencia y en comité quincenal; si falla la hipótesis, conservar la valoración inicial y activar 8.F.1. \\
R-021 & 15,00\% / 2 / 0,30 & 2 / 0,10; 0,03 (Baja) & L=13,33 HH; p1L=2,00 HH. & No se reduce p0: Ventana y recorrido de prueba confirmados por Operaciones/transportista antes de salir. Revisar al verificar evidencia y en comité quincenal; si falla la hipótesis, conservar la valoración inicial y activar 8.F.1. \\
R-022 & 10,00\% / 1 / 0,10 & 2 / 0,10; 0,01 (Baja) & L=13,33 HH; p1L=1,33 HH. & No se reduce p0: Autorización de recorrido y protocolo de seguridad del CLIENTE/transportista. Revisar al verificar evidencia y en comité quincenal; si falla la hipótesis, conservar la valoración inicial y activar 8.F.1. \\
R-023 & 10,00\% / 1 / 0,10 & 2 / 0,10; 0,01 (Baja) & L=17,33 HH; p1L=1,73 HH. & No se reduce p0: Custodia y seguridad del transporte corresponden al CLIENTE/transportista; evidencia de prueba preservada. Revisar al verificar evidencia y en comité quincenal; si falla la hipótesis, conservar la valoración inicial y activar 8.F.1. \\
R-024 & 30,00\% / 2 / 0,30 & 3 / 0,20; 0,06 (Baja) & L=27,00 HH; p1L=8,10 HH. & Hipótesis p1=30,00\% por resolución temprana de la dependencia: Identidad, autorización, capacitación y vínculo persona--vehículo--viaje verificados en andén. Revisar al verificar evidencia y en comité quincenal; si falla la hipótesis, conservar la valoración inicial y activar 8.F.1. \\
R-025 & 30,00\% / 2 / 0,30 & 3 / 0,20; 0,06 (Baja) & L=73,33 HH; p1L=22,00 HH. & Hipótesis p1=30,00\% por resolución temprana de la dependencia: Solicitud y análisis de impacto aprobados antes de iniciar trabajo fuera de alcance. Revisar al verificar evidencia y en comité quincenal; si falla la hipótesis, conservar la valoración inicial y activar 8.F.1. \\
R-026 & 20,00\% / 2 / 0,30 & 4 / 0,40; 0,12 (Media) & L=69,33 HH; p1L=13,87 HH. & No se reduce p0: Seguimiento de cambios aplicables con asesoría competente; decisión contractual documentada. Revisar al verificar evidencia y en comité quincenal; si falla la hipótesis, conservar la valoración inicial y activar 8.F.1. \\
R-027 & 25,00\% / 2 / 0,30 & 3 / 0,20; 0,06 (Baja) & L=49,33 HH; p1L=12,33 HH. & Hipótesis p1=25,00\% por resolución temprana de la dependencia: Reglas de costo de servir aprobadas al mes 20 y tres liquidaciones en sombra reproducidas en 21--23. Revisar al verificar evidencia y en comité quincenal; si falla la hipótesis, conservar la valoración inicial y activar 8.F.1. \\
R-028 & 25,00\% / 2 / 0,30 & 4 / 0,40; 0,12 (Media) & L=56,00 HH; p1L=14,00 HH. & Hipótesis p1=25,00\% por resolución temprana de la dependencia: Acuerdo 5.4.2 con finalidades, accesos y controles comprendidos por Operaciones/TI/sindicato. Revisar al verificar evidencia y en comité quincenal; si falla la hipótesis, conservar la valoración inicial y activar 8.F.1. \\
R-029 & 35,00\% / 3 / 0,50 & 4 / 0,40; 0,20 (Alta) & L=176,00 HH; p1L=61,60 HH. & Hipótesis p1=35,00\% por resolución temprana de la dependencia: Altas por turno, tutor, evaluación antes de autonomía y aprendizaje dentro de dos horas según 7.1.2. Plan adicional en 8.F.1; mantener bloqueo/escalamiento y revisar diariamente hasta cerrar evidencia. \\
R-030 & 35,00\% / 3 / 0,50 & 4 / 0,40; 0,20 (Alta) & L=116,00 HH; p1L=40,60 HH. & Hipótesis p1=35,00\% por resolución temprana de la dependencia: Contrato de extracción e intercambio confirmado con TI; ensayo temprano sin modificar ERP y sustitución progresiva de WMS. Plan adicional en 8.F.1; mantener bloqueo/escalamiento y revisar diariamente hasta cerrar evidencia. \\
R-031 & 35,00\% / 3 / 0,50 & 4 / 0,40; 0,20 (Alta) & L=57,33 HH; p1L=20,07 HH. & Hipótesis p1=35,00\% por resolución temprana de la dependencia: Disponibilidad y continuidad acordadas antes de septiembre; seguimiento diario de las seis condiciones de cierre. Plan adicional en 8.F.1; mantener bloqueo/escalamiento y revisar diariamente hasta cerrar evidencia. \\
R-032 & 25,00\% / 2 / 0,30 & 4 / 0,40; 0,12 (Media) & L=86,67 HH; p1L=21,67 HH. & Hipótesis p1=25,00\% por resolución temprana de la dependencia: Inventario de versiones y soporte, revisión trimestral y compatibilidad probada antes del límite verificado. Revisar al verificar evidencia y en comité quincenal; si falla la hipótesis, conservar la valoración inicial y activar 8.F.1. \\
R-033 & 20,00\% / 2 / 0,30 & 4 / 0,40; 0,12 (Media) & L=109,33 HH; p1L=21,87 HH. & Hipótesis p1=20,00\% por resolución temprana de la dependencia: Cuentas, claves, exportación y restauración utilizable verificadas en recepción y revisión anual de salida. Revisar al verificar evidencia y en comité quincenal; si falla la hipótesis, conservar la valoración inicial y activar 8.F.1. \\
R-034 & 25,00\% / 2 / 0,30 & 4 / 0,40; 0,12 (Media) & L=89,33 HH; p1L=22,33 HH. & Hipótesis p1=25,00\% por resolución temprana de la dependencia: Procedimiento ERP con TI/Finanzas probado en AL-DTE-01/AL-DR-01 antes de H5/H10. Revisar al verificar evidencia y en comité quincenal; si falla la hipótesis, conservar la valoración inicial y activar 8.F.1. \\
R-035 & 25,00\% / 2 / 0,30 & 4 / 0,40; 0,12 (Media) & L=57,33 HH; p1L=14,33 HH. & Hipótesis p1=25,00\% por resolución temprana de la dependencia: Demostración práctica en andén e incorporación acordada con transportista y conductor habilitado. Revisar al verificar evidencia y en comité quincenal; si falla la hipótesis, conservar la valoración inicial y activar 8.F.1. \\
\end{tablalafrox}
\end{landscape}
\endgroup

\begingroup
\renewcommand{\thetable}{8.F.3}
\begin{landscape}
\setlength{\textwidth}{\linewidth}
\renewcommand{\arraystretch}{1.08}
\begin{tablalafrox}{Costo-beneficio y sensibilidad de las 15 respuestas altas/críticas iniciales (HH). Fuente: modelo explícito de 8.F.}{tab:cierre-cb}{G{0.06} F{0.22} F{0.22} F{0.25} Y}{\cab{ID} & \cab{Entradas acumuladas} & \cab{Resultado esperado} & \cab{Sensibilidad / umbral} & \cab{Decisión y justificación}}
R-001 & A=8,00; C=16,00; H0=104,00; p0=45,00\%; p1=25,00\%. & B=24,80; T=12,00; B/T=2,07; neto=12,80 HH. & Sin mejora p1=p0: B/T=0,47; neto=-8,00 HH. Equilibrio neto: p1<=37,31\%. & Eficiencia de esfuerzo favorable bajo la hipótesis p1; mantener el tratamiento por obligación y no usar B/T como autorización de aceptación. \\
R-002 & A=12,00; C=24,00; H0=194,67; p0=60,00\%; p1=60,00\%. & B=14,40; T=26,40; B/T=0,55; neto=-12,00 HH. & Sin mejora p1=p0: B/T=0,55; neto=-12,00 HH. Equilibrio neto: p1<=53,84\%. & Prioridad obligatoria: proteger trazabilidad, condiciones de aceptación y continuidad; se mantiene aun con retorno HH menor que uno. Validar capacidad y elevar su esfuerzo antes del hito. \\
R-004 & A=8,00; C=20,00; H0=110,67; p0=55,00\%; p1=30,00\%. & B=33,67; T=14,00; B/T=2,40; neto=19,67 HH. & Sin mejora p1=p0: B/T=0,58; neto=-8,00 HH. Equilibrio neto: p1<=47,77\%. & Eficiencia de esfuerzo favorable bajo la hipótesis p1; mantener el tratamiento por obligación y no usar B/T como autorización de aceptación. \\
R-005 & A=6,00; C=16,00; H0=96,00; p0=50,00\%; p1=30,00\%. & B=24,00; T=10,80; B/T=2,22; neto=13,20 HH. & Sin mejora p1=p0: B/T=0,57; neto=-6,00 HH. Equilibrio neto: p1<=43,75\%. & Eficiencia de esfuerzo favorable bajo la hipótesis p1; mantener el tratamiento por obligación y no usar B/T como autorización de aceptación. \\
R-006 & A=8,00; C=16,00; H0=101,33; p0=55,00\%; p1=35,00\%. & B=25,87; T=13,60; B/T=1,90; neto=12,27 HH. & Sin mejora p1=p0: B/T=0,52; neto=-8,00 HH. Equilibrio neto: p1<=47,11\%. & Eficiencia de esfuerzo favorable bajo la hipótesis p1; mantener el tratamiento por obligación y no usar B/T como autorización de aceptación. \\
R-009 & A=6,00; C=12,00; H0=77,33; p0=35,00\%; p1=20,00\%. & B=14,00; T=8,40; B/T=1,67; neto=5,60 HH. & Sin mejora p1=p0: B/T=0,41; neto=-6,00 HH. Equilibrio neto: p1<=27,24\%. & Eficiencia de esfuerzo favorable bajo la hipótesis p1; mantener el tratamiento por obligación y no usar B/T como autorización de aceptación. \\
R-011 & A=4,00; C=8,00; H0=43,33; p0=35,00\%; p1=25,00\%. & B=6,33; T=6,00; B/T=1,06; neto=0,33 HH. & Sin mejora p1=p0: B/T=0,41; neto=-4,00 HH. Equilibrio neto: p1<=25,77\%. & Eficiencia de esfuerzo favorable bajo la hipótesis p1; mantener el tratamiento por obligación y no usar B/T como autorización de aceptación. \\
R-013 & A=8,00; C=20,00; H0=128,00; p0=55,00\%; p1=30,00\%. & B=38,00; T=14,00; B/T=2,71; neto=24,00 HH. & Sin mejora p1=p0: B/T=0,58; neto=-8,00 HH. Equilibrio neto: p1<=48,75\%. & Eficiencia de esfuerzo favorable bajo la hipótesis p1; mantener el tratamiento por obligación y no usar B/T como autorización de aceptación. \\
R-016 & A=4,00; C=12,00; H0=85,33; p0=45,00\%; p1=25,00\%. & B=20,07; T=7,00; B/T=2,87; neto=13,07 HH. & Sin mejora p1=p0: B/T=0,57; neto=-4,00 HH. Equilibrio neto: p1<=40,31\%. & Eficiencia de esfuerzo favorable bajo la hipótesis p1; mantener el tratamiento por obligación y no usar B/T como autorización de aceptación. \\
R-028 & A=6,00; C=12,00; H0=68,00; p0=45,00\%; p1=25,00\%. & B=16,60; T=9,00; B/T=1,84; neto=7,60 HH. & Sin mejora p1=p0: B/T=0,47; neto=-6,00 HH. Equilibrio neto: p1<=36,18\%. & Eficiencia de esfuerzo favorable bajo la hipótesis p1; mantener el tratamiento por obligación y no usar B/T como autorización de aceptación. \\
R-029 & A=12,00; C=32,00; H0=208,00; p0=55,00\%; p1=35,00\%. & B=52,80; T=23,20; B/T=2,28; neto=29,60 HH. & Sin mejora p1=p0: B/T=0,59; neto=-12,00 HH. Equilibrio neto: p1<=49,23\%. & Eficiencia de esfuerzo favorable bajo la hipótesis p1; mantener el tratamiento por obligación y no usar B/T como autorización de aceptación. \\
R-030 & A=12,00; C=24,00; H0=140,00; p0=50,00\%; p1=35,00\%. & B=29,40; T=20,40; B/T=1,44; neto=9,00 HH. & Sin mejora p1=p0: B/T=0,50; neto=-12,00 HH. Equilibrio neto: p1<=41,43\%. & Eficiencia de esfuerzo favorable bajo la hipótesis p1; mantener el tratamiento por obligación y no usar B/T como autorización de aceptación. \\
R-031 & A=8,00; C=16,00; H0=73,33; p0=45,00\%; p1=35,00\%. & B=12,93; T=13,60; B/T=0,95; neto=-0,67 HH. & Sin mejora p1=p0: B/T=0,47; neto=-8,00 HH. Equilibrio neto: p1<=34,09\%. & Prioridad obligatoria: proteger trazabilidad, condiciones de aceptación y continuidad; se mantiene aun con retorno HH menor que uno. Validar capacidad y elevar su esfuerzo antes del hito. \\
R-032 & A=6,00; C=20,00; H0=106,67; p0=40,00\%; p1=25,00\%. & B=21,00; T=11,00; B/T=1,91; neto=10,00 HH. & Sin mejora p1=p0: B/T=0,57; neto=-6,00 HH. Equilibrio neto: p1<=34,38\%. & Eficiencia de esfuerzo favorable bajo la hipótesis p1; mantener el tratamiento por obligación y no usar B/T como autorización de aceptación. \\
R-035 & A=6,00; C=16,00; H0=73,33; p0=40,00\%; p1=25,00\%. & B=15,00; T=10,00; B/T=1,50; neto=5,00 HH. & Sin mejora p1=p0: B/T=0,52; neto=-6,00 HH. Equilibrio neto: p1<=31,82\%. & Eficiencia de esfuerzo favorable bajo la hipótesis p1; mantener el tratamiento por obligación y no usar B/T como autorización de aceptación. \\
\end{tablalafrox}
\end{landscape}
\endgroup
